The denominator is $(x^3 - \dfrac{1}{2})^2 + \dfrac{3}{4} > 0$. For $x\ne0$, set $t = x + \dfrac{1}{x}$, so $t\le-2$ or $t\ge2$, and $f = \dfrac{(t-1)^3}{t^3-3t-1}$. Differentiation with respect to $t$ gives $\dfrac{3 (t-1)^2 (t^2-2t-2)}{(t^3-3t-1)^2}$. On $t\le-2$ this is positive, and $f$ increases from limit 1 to $f(-2)=9$. On $t\ge2$, it decreases until $t=1+\sqrt3$ and then increases to limit 1. The minimum is $\dfrac{\sqrt{3}}{1+\sqrt3}$; it is attained at the two positive roots of $x^2-(1+\sqrt3)x+1=0$. Also $f(0)=1$. Hence the minimum is $\boxed{\dfrac{\sqrt{3}}{1+\sqrt3}}$ and the maximum is $\boxed9$, attained at $x=-1$.
